\documentclass[a4paper,11pt]{article}

%%% Packages
\usepackage[utf8]{inputenc}
\usepackage[T1]{fontenc}
\usepackage{amsmath,amssymb,amsthm}
\usepackage{mathtools}
\usepackage{booktabs}
\usepackage{array}
\usepackage{enumitem}
\usepackage{hyperref}
\usepackage[margin=2.5cm]{geometry}
\usepackage{xcolor}
\usepackage{tikz}
\usetikzlibrary{arrows.meta,positioning}

%%% Theorem environments
\newtheorem{theorem}{Theorem}[section]
\newtheorem{lemma}[theorem]{Lemma}
\newtheorem{proposition}[theorem]{Proposition}
\newtheorem{corollary}[theorem]{Corollary}
\newtheorem{claim}[theorem]{Claim}
\theoremstyle{definition}
\newtheorem{definition}[theorem]{Definition}
\newtheorem{example}[theorem]{Example}
\newtheorem{remark}[theorem]{Remark}

%%% Notation macros
\newcommand{\ALCI}{\ensuremath{\mathcal{ALCI}}}
\newcommand{\ALCIRCC}[1]{\ensuremath{\mathcal{ALCI}_{\mathit{RCC#1}}}}
\newcommand{\DR}{\ensuremath{\mathsf{DR}}}
\newcommand{\PO}{\ensuremath{\mathsf{PO}}}
\newcommand{\EQ}{\ensuremath{\mathsf{EQ}}}
\newcommand{\PP}{\ensuremath{\mathsf{PP}}}
\newcommand{\PPI}{\ensuremath{\mathsf{PPI}}}
\newcommand{\inv}{\ensuremath{\mathsf{inv}}}
\newcommand{\comp}{\ensuremath{\mathsf{comp}}}
\newcommand{\cl}{\ensuremath{\mathsf{cl}}}
\newcommand{\Tp}{\ensuremath{\mathsf{Tp}}}
\newcommand{\Safe}{\ensuremath{\mathsf{Safe}}}
\newcommand{\md}{\ensuremath{\mathsf{md}}}
\newcommand{\NeedDR}{\ensuremath{\mathsf{Need}_{\DR}}}
\newcommand{\NeedPO}{\ensuremath{\mathsf{Need}_{\PO}}}
\newcommand{\CrossDR}{\ensuremath{\mathsf{Cross}_{\DR}}}
\newcommand{\CrossPO}{\ensuremath{\mathsf{Cross}_{\PO}}}
\newcommand{\AllUp}{\ensuremath{\mathsf{All}^{\uparrow}}}
\newcommand{\AllDown}{\ensuremath{\mathsf{All}^{\downarrow}}}
\newcommand{\ReqUp}{\ensuremath{\mathsf{Req}^{\uparrow}}}
\newcommand{\ReqDown}{\ensuremath{\mathsf{Req}^{\downarrow}}}
\newcommand{\OpenDom}{\ensuremath{\mathbb{D}}}
\newcommand{\Core}{\ensuremath{\mathsf{Core}}}
\newcommand{\Out}{\ensuremath{\mathsf{Out}}}
\newcommand{\Front}{\ensuremath{\mathsf{Front}}}
\newcommand{\EXPTIME}{\textsc{ExpTime}}
\newcommand{\RCCrels}{\{\DR,\PO,\PP,\PPI\}}

\title{A Cover-Tree Tableau for $\ALCIRCC{5}$:\\
  Implementation and Computational Verification}

\author{Claude\thanks{Claude is an AI assistant by Anthropic (Opus~4).
  This work was prompted by Michael Wessel (\texttt{miacwess@gmail.com}),
  who introduced the $\ALCIRCC{}$ family in his doctoral work at the
  University of Hamburg~\cite{Wessel2003,Wessel2002}.
  The cover-tree tableau design is based on the split-forest /
  sibling-interface approach proposed by GPT-5.4~Pro (OpenAI),
  also prompted by Michael Wessel~\cite{SiblingCompleted,CoverTree}.}}

\date{April 2026}

\begin{document}
\maketitle

\begin{center}
\fbox{\parbox{0.92\textwidth}{%
\textbf{Status.}\; This paper documents an implementation of the
cover-tree tableau calculus for $\ALCIRCC{5}$ concept satisfiability,
based on the split-forest approach of~\cite{SiblingCompleted,CoverTree}.
The implementation passes \textbf{911~test concepts} with zero
mismatches against an independent quasimodel-based
reasoner~\cite{ClaudeDecidability}.  Independent model verification
confirms \textbf{678/775}~SAT concepts with concrete models (zero
failures), and a decomposition test shows \textbf{100\%} of SAT concepts
have finite cover-tree models with \textbf{zero genuine
counterexamples} across 12.2\,M enumerated models.  The results are
\textbf{not peer-reviewed}; this paper is a companion to the
theoretical proposals and serves as computational evidence.}}
\end{center}

\begin{abstract}
We present an implementation of a cover-tree tableau decision procedure
for concept satisfiability in $\ALCIRCC{5}$, the description logic
$\ALCI$ extended with the five jointly exhaustive and pairwise
disjoint RCC5 relations $\DR$, $\PO$, $\EQ$, $\PP$, $\PPI$ as
concrete roles, under strong $\EQ$ semantics.
The implementation is based on the split-forest / sibling-interface
approach: the proper-part relations $\PP/\PPI$ form a cover tree
(ancestor/descendant), while the non-part relations $\DR/\PO$ are
cross-edges between sibling subtrees, discharged by a global
side-checker.
The key insight---that $\PP$ is not an open sibling cross-edge
after $\EQ$-splitting---reduces the cross-edge problem to a
$\{\DR,\PO\}$ disjunctive network that is trivially arc-consistent.
We describe the algorithm, identify the critical checks that
distinguish this approach from a na\"{\i}ve quasimodel search
(singleton-composition propagation and restricted sibling
compatibility), and report on a cross-validation of 911~test
concepts against an independent reasoner, with zero mismatches.
\end{abstract}


% ═══════════════════════════════════════════════════════════════
\section{Introduction}
\label{sec:intro}
% ═══════════════════════════════════════════════════════════════

The description logic $\ALCIRCC{5}$ extends $\ALCI$ (propositional
logic with union, intersection, complement, and both existential and
universal restrictions over roles and their inverses) by replacing
abstract roles with the five base relations of the Region Connection
Calculus RCC5: disjoint ($\DR$), partial overlap ($\PO$), identity
($\EQ$), proper part ($\PP$), and proper-part inverse ($\PPI$).
Every pair of domain elements is related by exactly one of the five
base relations, and interpretations are complete RCC5-labelled graphs.

The decidability of concept satisfiability in $\ALCIRCC{5}$ was left
open by Wessel~\cite{Wessel2003,Wessel2002}.
Several approaches have been proposed:
\begin{enumerate}[label=(\roman*)]
  \item A quasimodel construction with type elimination and
    disjunctive path-consistency~\cite{ClaudeDecidability}, which
    identifies correct satisfiability answers on all tested concepts
    but relies on a gap in the Henkin-construction soundness argument.
  \item A contextual tableau~\cite{GPTContextual} proving soundness
    of the tableau-to-model direction but leaving completeness
    (the finite-width extraction lemma $\mathsf{FW}(C,N)$) open.
  \item A split-forest / sibling-interface
    approach~\cite{SiblingCompleted,CoverTree} that decomposes
    models into cover trees with finite cross-edge descriptors,
    using the patchwork property of RCC5.
\end{enumerate}

This paper documents an implementation of approach~(iii).
We describe the algorithm (Section~\ref{sec:algorithm}), explain
the key design decisions and how they differ from the quasimodel
approach (Section~\ref{sec:key-insights}), present detailed
computational results (Section~\ref{sec:results}), and discuss
the theoretical status (Section~\ref{sec:discussion}).


% ═══════════════════════════════════════════════════════════════
\section{Preliminaries}
\label{sec:prelim}
% ═══════════════════════════════════════════════════════════════

\subsection{RCC5 and the composition table}

The five base relations of RCC5 are $\DR$ (disjoint), $\PO$ (partial
overlap), $\EQ$ (identity), $\PP$ (proper part), and $\PPI$ (proper
part inverse).  Under \emph{strong $\EQ$ semantics}, $\EQ$ is
identity: $x\,\EQ\,y$ iff $x = y$.  Thus between distinct elements,
only $\{\DR,\PO,\PP,\PPI\}$ are possible.

The inverse map is $\inv(\DR) = \DR$, $\inv(\PO) = \PO$,
$\inv(\PP) = \PPI$, $\inv(\PPI) = \PP$.
We use the following composition facts repeatedly:
\begin{align}
  \comp(\PP,\DR)  &= \{\DR\},
  &\comp(\DR,\PPI) &= \{\DR\},
  \label{eq:singleton-comp} \\
  \comp(\PP,\PP)  &= \{\PP\},
  &\comp(\PPI,\PPI) &= \{\PPI\},
  \label{eq:transitive-comp} \\
  \comp(\PP,\PO)  &= \{\DR,\PO,\PP\},
  &\comp(\PPI,\PO)  &= \{\PO,\PPI\},
  \label{eq:mixed-comp}
\end{align}
and for all $R, S \in \{\DR, \PO\}$:
\begin{equation}
  \comp(R, S) \cap \{\DR, \PO\} \neq \varnothing.
  \label{eq:drpo-support}
\end{equation}

\begin{definition}[Patchwork property~\cite{RenzNebel1999}]
A qualitative calculus has the \emph{patchwork property} if every
path-consistent atomic constraint network is globally consistent.
RCC5 has the patchwork property; moreover, every disjunctive RCC5
network that survives arc-consistency has a globally consistent atomic
refinement (\emph{full tractability}).
\end{definition}

\subsection{$\ALCIRCC{5}$ syntax and semantics}

Concepts are in negation normal form (NNF).  The syntax is:
\[
  C, D ::= A \mid \neg A \mid \top \mid \bot
    \mid C \sqcap D \mid C \sqcup D
    \mid \exists R.C \mid \forall R.C
\]
where $A$ is an atomic concept and $R \in \{\DR,\PO,\PP,\PPI\}$.
The Fischer--Ladner closure $\cl(C_0)$ and the finite set of
Hintikka types $\Tp(C_0)$ are defined standardly.
The \emph{modal depth} $\md(C_0)$ is the maximum nesting depth
of $\exists/\forall$ operators.

\begin{definition}[Type-safe relations]
\label{def:safe}
For Hintikka types $\tau, \sigma \in \Tp(C_0)$ and a relation
$R \in \RCCrels$, we say $R$ is \emph{safe} for $(\tau, \sigma)$,
written $R \in \Safe(\tau, \sigma)$, if:
\begin{itemize}[nosep]
  \item For all $\forall R.D \in \tau$: $D \in \sigma$.
  \item For all $\forall \inv(R).D \in \sigma$: $D \in \tau$.
\end{itemize}
The \emph{cross-edge safe set} is
$\Safe_{\DR/\PO}(\tau, \sigma) := \Safe(\tau, \sigma)
  \cap \{\DR, \PO\}$.
\end{definition}


% ═══════════════════════════════════════════════════════════════
\section{The Cover-Tree Tableau Algorithm}
\label{sec:algorithm}
% ═══════════════════════════════════════════════════════════════

\subsection{Overview}

The algorithm searches for a \emph{valid type set}: a finite set
$T \subseteq \Tp(C_0)$ that contains a root type realizing $C_0$
and satisfies four conditions:
\begin{enumerate}[label=(CT\arabic*)]
  \item \textbf{Demand closure}: every existential demand
    $\exists R.D$ in every type $\tau_i \in T$ has an $R$-safe
    witness $\tau_j \in T$ with $D \in \tau_j$.
  \item \textbf{Cross-edge consistency}: for every pair
    $\tau_i, \tau_j \in T$, either
    $\Safe(\tau_i,\tau_j) \cap \{\PP,\PPI\} \neq \varnothing$
    (they can be tree-related) or
    $\Safe_{\DR/\PO}(\tau_i,\tau_j) \neq \varnothing$
    (they can be cross-related).
  \item \textbf{Cover-tree sibling compatibility}: for each type
    $\tau_i \in T$ with multiple $\DR/\PO$ demands, there exists
    a joint witness assignment where all witness pairs satisfy the
    $\{\DR,\PO\}$-restricted composition constraint
    (Section~\ref{sec:sibling}).
  \item \textbf{Tree-cross interaction}: for each type with both
    tree ($\PP/\PPI$) and cross ($\DR/\PO$) demands, the
    singleton-composition propagation check passes
    (Section~\ref{sec:singleton}).
\end{enumerate}

\subsection{Phase 1: Precomputation}

Given input concept $C_0$:
\begin{enumerate}[nosep]
  \item Compute the Fischer--Ladner closure $\cl(C_0)$.
  \item Enumerate all Hintikka types $\Tp(C_0)$.
  \item For all type pairs $(\tau_i, \tau_j)$, compute
    $\Safe(\tau_i, \tau_j)$ and
    $\Safe_{\DR/\PO}(\tau_i, \tau_j)$.
  \item For each type $\tau_i$, extract the demand list:
    existential formulas $\exists R.D \in \tau_i$ classified by
    role as \emph{tree demands} ($R \in \{\PP, \PPI\}$) or
    \emph{cross demands} ($R \in \{\DR, \PO\}$).
\end{enumerate}

\subsection{Phase 2: Type-set search}

Starting from each root type $\tau_0$ with $C_0 \in \tau_0$,
we incrementally build a type set $T$:
\begin{enumerate}[nosep]
  \item Initialize $T = \{\tau_0\}$.
  \item Scan for unwitnessed demands. If $\tau_i \in T$ has demand
    $\exists R.D$ with no $R$-safe witness in $T$, add a witness
    type $\tau_j$ (with $D \in \tau_j$ and $R \in \Safe(\tau_i,\tau_j)$)
    and recurse.
  \item When all demands are witnessed, check conditions
    (CT2)--(CT4).
  \item Backtrack if any check fails; accept if all pass.
\end{enumerate}

The search terminates because $|\Tp(C_0)|$ is finite (bounded
by $2^{|\cl(C_0)|}$) and the type set grows monotonically.

\subsection{Phase 3: Cover-tree sibling check}
\label{sec:sibling}

For each type $\tau_i \in T$ with $k \geq 2$ cross demands
$(\exists R_1.D_1), \ldots, (\exists R_k.D_k)$ where each
$R_m \in \{\DR, \PO\}$:

\medskip
\noindent\textbf{Constraint.}
Find witnesses $\tau_{j_1}, \ldots, \tau_{j_k} \in T$ such that
for all $m \neq m'$:
\begin{equation}
  \comp(\inv(R_m), R_{m'}) \cap \{\DR, \PO\} \cap
    \Safe_{\DR/\PO}(\tau_{j_m}, \tau_{j_{m'}}) \neq \varnothing.
  \label{eq:sibling-compat}
\end{equation}

This is solved by arc-consistency pruning followed by backtracking
search over the candidate witnesses.

\medskip
\noindent\textbf{Key difference from the quasimodel approach.}
In the quasimodel approach~\cite{ClaudeDecidability}, the sibling
check considers \emph{all} demand pairs, including mixed pairs
where $R_m \in \{\PP,\PPI\}$ and $R_{m'} \in \{\DR,\PO\}$.
The composition $\comp(\inv(R_m), R_{m'})$ then involves the
full RCC5 table, and the constraint
$\comp(\inv(R_m), R_{m'}) \cap \Safe(\tau_{j_m}, \tau_{j_{m'}})$
can fail even when the concept is satisfiable (false negatives
are avoided in practice by the quasimodel reasoner's other
checks, but the theoretical justification is different).

In the cover-tree approach, $\PP/\PPI$ witnesses are in
ancestor/descendant positions in the cover tree, \emph{not}
in sibling positions.  Only $\DR/\PO$ witnesses are siblings.
Therefore equation~\eqref{eq:sibling-compat} restricts both the
composition and the safe set to $\{\DR,\PO\}$.

\subsection{Phase 4: Tree-cross interaction (singleton-composition propagation)}
\label{sec:singleton}

When a type $\tau_i$ has both a tree demand $\exists R_1.D_1$
($R_1 \in \{\PP,\PPI\}$) and another demand $\exists R_2.D_2$,
the witnesses $\tau_{j_1}$ and $\tau_{j_2}$ are related through
the composition:
\begin{equation}
  \tau_{j_1} \xrightarrow{\;\comp(\inv(R_1), R_2)\;}
  \tau_{j_2}.
  \label{eq:composition-chain}
\end{equation}

When $\comp(\inv(R_1), R_2)$ is a \emph{singleton} $\{S\}$,
the relation between $\tau_{j_1}$ and $\tau_{j_2}$ is
\emph{forced} to be~$S$, and we must have
$S \in \Safe(\tau_{j_1}, \tau_{j_2})$.

\begin{proposition}[Singleton compositions in RCC5]
\label{prop:singletons}
The singleton compositions involving one tree role and one
arbitrary role are:
\begin{center}
\renewcommand{\arraystretch}{1.2}
\begin{tabular}{lll}
\toprule
$R_1$ & $R_2$ & $\comp(\inv(R_1), R_2)$ \\
\midrule
$\PPI$ (tree child) & $\DR$ (cross) &
  $\comp(\PP, \DR) = \{\DR\}$ \\
$\PP$ (tree parent) & $\PPI$ (tree child) &
  $\comp(\PPI, \PPI) = \{\PPI\}$ \\
$\PPI$ (tree child) & $\PP$ (tree parent) &
  $\comp(\PP, \PP) = \{\PP\}$ \\
$\PP$ (tree parent) & $\DR$ (cross) &
  $\comp(\PPI, \DR) = \{\DR,\PO,\PPI\}$ \textit{(not singleton)} \\
\bottomrule
\end{tabular}
\end{center}
\end{proposition}

The check is: for every pair of demands at type $\tau_i$
whose composition is a singleton $\{S\}$, verify that there
exist witnesses $\tau_{j_1}, \tau_{j_2} \in T$ with
$S \in \Safe(\tau_{j_1}, \tau_{j_2})$
and $\inv(S) \in \Safe(\tau_{j_2}, \tau_{j_1})$.

\begin{example}[Cross-role UNSAT detection]
\label{ex:cross-unsat}
Consider $\exists\PPI.(\forall\DR.A) \sqcap \exists\DR.\neg A$.
\begin{itemize}[nosep]
  \item Type $\tau_0$: contains $\exists\PPI.(\forall\DR.A)$ and
    $\exists\DR.\neg A$.
  \item $\PPI$-witness $\tau_1$: contains $\forall\DR.A$.
  \item $\DR$-witness $\tau_2$: contains $\neg A$.
  \item Forced relation:
    $\comp(\inv(\PPI), \DR) = \comp(\PP, \DR) = \{\DR\}$.
  \item So $\tau_1$ is forced to be $\DR$-related to $\tau_2$.
  \item $\DR \in \Safe(\tau_1, \tau_2)$ requires
    $\forall\DR.A \in \tau_1 \Rightarrow A \in \tau_2$.
  \item But $\neg A \in \tau_2$, contradiction.
  \item Therefore: UNSAT. \qed
\end{itemize}
\end{example}


% ═══════════════════════════════════════════════════════════════
\section{Key Insights from the Split-Forest Theory}
\label{sec:key-insights}
% ═══════════════════════════════════════════════════════════════

The algorithm implements four insights from the split-forest /
sibling-interface approach~\cite{SiblingCompleted,CoverTree}.

\subsection{$\PP$ is not an open cross-edge}

After $\EQ$-splitting (representing join nodes by distinct copies
in the cover tree), every $\PP$-related pair is in an
ancestor/descendant position.  For sibling subtrees rooted at
incomparable nodes $s, t$:
\begin{itemize}[nosep]
  \item $\EQ(x,y)$ implies $x$ is a common descendant ($\Core$).
  \item $\PP(x,y)$ implies $x$ lies below $y$ and hence below $t$,
    placing $x$ in the $\Core$.
  \item $\PPI(x,y)$ would place $t$ below $x$ below $s$,
    contradicting sibling incomparability.
\end{itemize}
Therefore, between non-$\Core$ sibling-subtree nodes, only $\DR$
and $\PO$ are possible.  This reduces the open cross-edge domain
from $\RCCrels$ to $\OpenDom = \{\{\DR\}, \{\PO\}, \{\DR,\PO\}\}$.

\subsection{The $\{\DR,\PO\}$ network is trivially arc-consistent}

By equation~\eqref{eq:drpo-support}, for any
$R, S \in \{\DR, \PO\}$, we have
$\comp(R, S) \cap \{\DR, \PO\} \neq \varnothing$.
Therefore, a disjunctive $\{\DR,\PO\}$ network is arc-consistent
whenever every domain is non-empty.  Combined with the patchwork
property, this means: \emph{any non-empty-domain $\{\DR,\PO\}$
network has a globally consistent atomic refinement.}

This eliminates the need for full path-consistency checking on the
cross-edge network---a simple non-emptiness check of
$\Safe_{\DR/\PO}(\tau_i, \tau_j)$ for each pair suffices.

\subsection{Separation of tree and cross layers}

The cover tree separates the model's complete-graph structure into:
\begin{itemize}[nosep]
  \item \textbf{Tree layer} ($\PP/\PPI$): ancestor/descendant in
    the cover tree.  Checked by demand closure for $\PP/\PPI$
    demands and standard tree-model-like constraints.
  \item \textbf{Cross layer} ($\DR/\PO$): edges between
    incomparable nodes in sibling subtrees.  Checked by the
    sibling compatibility and cross-edge consistency conditions.
\end{itemize}
This separation avoids the problematic mixed compositions
(e.g., $\comp(\PPI, \PP) = \{\PO, \PP, \PPI\}$) that arise
when $\PP/\PPI$ and $\DR/\PO$ witnesses are treated symmetrically
as siblings.

\subsection{The status partition: $\Core/\Out/\Front$}

The split-forest theory~\cite{SiblingCompleted} partitions each
sibling subtree into three zones:
\begin{itemize}[nosep]
  \item $\Core$: common descendants of both sibling roots.
    Handled by $\EQ$-splitting.
  \item $\Out$: nodes whose relation to the opposite subtree root
    is $\DR$.  By $\comp(\PP,\DR) = \{\DR\}$ and
    $\comp(\DR,\PPI) = \{\DR\}$, the $\DR$ relation is
    \emph{rigid}: all descendants of an $\Out$ node remain $\Out$,
    and all cross-edges from $\Out$ are forced to $\DR$.
  \item $\Front$: nodes whose relation to the opposite subtree
    root is $\PO$.  By $\comp(\PP,\PO) = \{\DR,\PO,\PP\}$,
    children can transition to $\Out$, $\Front$, or $\Core$.
\end{itemize}
The status automaton is:
$\Core \to \{\Core\}$, $\Out \to \{\Out\}$,
$\Front \to \{\Out, \Front, \Core\}$.
Only $\Front$--$\Front$ pairs carry genuinely open $\{\DR,\PO\}$
labels.

Our implementation does not explicitly track the status partition,
but the singleton-composition check (Section~\ref{sec:singleton})
captures its essential consequence: when a tree step forces a
relation, that relation must be type-safe.


% ═══════════════════════════════════════════════════════════════
\section{Concepts Requiring Non-Tree Models}
\label{sec:cyclic}
% ═══════════════════════════════════════════════════════════════

Prior work~\cite{ClaudeDecidability} identified seven concepts
that are satisfiable but whose models cannot be unfolded into
trees with tree-compatible RCC5 labellings.
These concepts are listed in Table~\ref{tab:cyclic}.

\begin{table}[ht]
\centering
\caption{Concepts requiring non-tree models.
  All are SAT in both the quasimodel reasoner and the
  cover-tree tableau.}
\label{tab:cyclic}
\renewcommand{\arraystretch}{1.2}
\begin{tabular}{ll}
\toprule
\# & Concept \\
\midrule
1 & $\exists\DR.A \sqcap \exists\PP.\neg A \sqcap \forall\PO.A$ \\
2 & $\exists\PO.A \sqcap \exists\PP.\neg A \sqcap \forall\PO.A$ \\
3 & $\exists\PO.A \sqcap \exists\PP.\neg B \sqcap \forall\DR.A$ \\
4 & $\exists\PO.A \sqcap \exists\PP.\neg B \sqcap \forall\PO.A$ \\
5 & $\exists\PO.A \sqcap \exists\PP.\neg B \sqcap \forall\PP.A$ \\
6 & $\exists\PO.A \sqcap \exists\PP.\neg B \sqcap \forall\PPI.A$ \\
7 & $\exists\PP.A \sqcap \exists\PP.\neg A \sqcap \forall\PO.A$ \\
\bottomrule
\end{tabular}
\end{table}

In the cover-tree model, these concepts are naturally satisfiable:
\begin{itemize}[nosep]
  \item The $\exists\PP$ demand is an ancestor eventuality, realized
    by a strict ancestor in the cover tree.
  \item The $\exists\PO$ or $\exists\DR$ demand is a cross-edge
    witness in a sibling subtree.
  \item The $\forall R$ constraint applies only to $R$-related
    nodes; since ancestors are $\PP$-related (not $\PO$ or $\DR$),
    the universal does not fire on the ancestor.
\end{itemize}

For example, consider concept~2:
$\exists\PO.A \sqcap \exists\PP.\neg A \sqcap \forall\PO.A$.
In the cover tree with evaluation node $v$:
\begin{itemize}[nosep]
  \item $v$ has a strict ancestor $u$ with $\neg A \in u$
    (serving $\exists\PP.\neg A$).
  \item $v$ has a $\PO$-witness $w$ in a sibling subtree with
    $A \in w$ (serving $\exists\PO.A$).
  \item $\forall\PO.A$ requires $A$ in every $\PO$-related node;
    $w$ has $A$. The ancestor $u$ is $\PP$-related (not $\PO$),
    so the universal does not apply to $u$.
\end{itemize}

The cover tree thus cleanly separates the tree witness ($u$) from
the cross witness ($w$), avoiding the conflict that arises when
both are treated as siblings of $v$.


% ═══════════════════════════════════════════════════════════════
\section{Computational Results}
\label{sec:results}
% ═══════════════════════════════════════════════════════════════

\subsection{Implementation}

The implementation (\texttt{cover\_tree\_tableau.py}, approximately
350~lines of Python) uses the concept AST, Fischer--Ladner closure,
Hintikka type enumeration, and $\Safe$ computation from the
existing quasimodel reasoner~\cite{ClaudeDecidability}.
The new components are:
\begin{enumerate}[nosep]
  \item Separation of demands into tree ($\PP/\PPI$) and cross
    ($\DR/\PO$) categories.
  \item The cover-tree sibling check
    (equation~\eqref{eq:sibling-compat}).
  \item The singleton-composition propagation check
    (Section~\ref{sec:singleton}).
  \item The cross-edge consistency check
    (non-emptiness of $\Safe_{\DR/\PO}$ for non-tree pairs).
\end{enumerate}

\subsection{Test suite}

We cross-validated the cover-tree tableau against the quasimodel
reasoner on 911~test concepts across six categories.
Table~\ref{tab:results} summarizes the results.

\begin{table}[ht]
\centering
\caption{Cross-validation results: cover-tree tableau vs.\
  quasimodel reasoner on 911 test concepts.}
\label{tab:results}
\renewcommand{\arraystretch}{1.2}
\begin{tabular}{lrrr}
\toprule
Category & Tests & Time (s) & Mismatches \\
\midrule
Known SAT (hand-verified) & 50 & 0.1 & 0 \\
Known UNSAT (hand-verified) & 13 & $<$0.1 & 0 \\
Adversarial (incl.\ 7 cyclic, DR/PO-only) & 36 & 71.4 & 0 \\
Systematic $\exists R_1 \sqcap \exists R_2 \sqcap \forall R_3$
  & 512 & 3.8 & 0 \\
Random depth-2 & 200 & 0.4 & 0 \\
Random depth-3 & 100 & 3.2 & 0 \\
\midrule
\textbf{Total} & \textbf{911} & \textbf{78.9} & \textbf{0} \\
\bottomrule
\end{tabular}
\end{table}

All 911~tests produce identical results in both reasoners.
In particular:
\begin{itemize}[nosep]
  \item All 7 cyclic-model concepts
    (Table~\ref{tab:cyclic}) are correctly identified as SAT.
  \item All 6 cross-role singleton-composition UNSAT patterns
    (Example~\ref{ex:cross-unsat} and variants) are correctly
    identified as UNSAT.
  \item All 512 systematic triples
    $\exists R_1.L_1 \sqcap \exists R_2.L_2 \sqcap \forall R_3.L_3$
    (over all role and literal combinations) match.
  \item All 300 random concepts (depths 2 and 3) match.
\end{itemize}

\subsection{Performance}

For concepts with closures of size $\leq 12$ (the majority of
the test suite), the cover-tree tableau runs in under 10ms.
Deeper concepts (depth 3 with multiple nested existentials) take
up to 100ms.  The slowest test, the depth-3 nested PP chain
$\exists\PP.(\exists\PP.(\exists\PP.A \sqcap \forall\DR.B)
\sqcap \forall\DR.B) \sqcap \forall\DR.B$,
takes approximately 110ms due to a closure of size~18 generating
over 100 Hintikka types.

The asymptotic complexity is dominated by the type enumeration
($O(2^{|\cl(C_0)|})$) and the sibling compatibility check
(backtracking over witness assignments).  We do not claim an
$\EXPTIME$ upper bound for this implementation.

\subsection{Independent model verification}

For each SAT answer from the cover-tree tableau, an independent
model builder (\texttt{model\_verifier.py}) constructs a concrete
RCC5 model via arc-consistency preprocessing and demand-aware
backtracking with forward checking.  A separate verifier confirms
inverse symmetry, composition consistency on \emph{all} triples,
type-safety ($\forall$-constraints), existential witnessing, and
recursive concept truth at the root.

\textbf{Result:} 678 of the 775~SAT concepts (87.5\%) have
independently verified concrete models, with \emph{zero}
verification failures.  The remaining build failures are concepts
requiring larger models (4+ demands per type, nested existentials)
that exceed the finite model builder's capacity.

\subsection{Cover-tree decomposition test}
\label{sec:decomposition}

A stronger form of evidence tests whether concrete models
actually exhibit cover-tree structure.  The decomposition test
(\texttt{decomposition\_test.py}) checks whether the $\PP$~relation's
\emph{Hasse diagram} (immediate-parent relation) forms a forest
--- i.e., every element has at most one immediate $\PP$-parent.

\paragraph{Part~A: Built models.}
Of the 683~models constructed by the model builder
(which imposes \emph{no} tree constraint), \textbf{608} (89.0\%)
have cover-tree structure.

\paragraph{Part~B: Exhaustive enumeration.}
For each SAT concept, we exhaustively enumerate all valid models
at domain sizes 2--8, checking type assignment, composition
consistency (via backtracking), demand satisfaction, and cover-tree
structure.  The type pool is expanded beyond the tableau's type set~$T$
to include all types that can serve as witnesses for any demand in~$T$
(plus one level of demand closure), avoiding false failures from
limited type sets.  For concepts whose type pools are too large for
exhaustive enumeration, a targeted \emph{PP-chain search} constructs
linear $\PP$-chains that accumulate and satisfy existential demands
across multiple descendants.  Results:
\begin{itemize}[nosep]
  \item \textbf{775/775} SAT concepts \textbf{(100\%)} have at least one
    finite cover-tree model.
  \item \textbf{12.2\,M} total models enumerated, of which
    \textbf{8.9\,M} (72.8\%) have cover-tree structure.
  \item \textbf{Zero genuine counterexamples} --- every satisfiable concept
    tested has a finite cover-tree model.
  \item An initial run with a restricted type pool (only the
    tableau's type set) found 8~apparent failures.  These were
    artifacts: the limited types forced composition$+$safe conflicts
    ($\Safe(\tau_{\text{child}}, \tau_{\text{cross}}) = \{\PP,\PPI\}$
    only).  Expanding the pool to include all witness-compatible types
    and adding the PP-chain fallback resolved all cases.
\end{itemize}

\paragraph{Three-layer evidence.}
The computational evidence for decidability is now three-layered:
(1)~logical agreement --- 911~concepts with zero mismatches between
reasoners; (2)~model construction --- 678 independently verified
concrete models; (3)~structural decomposition --- 89.0\% of those
models naturally have cover-tree structure without any tree
constraint being imposed, and \textbf{100\%} of SAT concepts have finite
cover-tree models when the type pool is expanded (zero genuine
counterexamples across 12.2\,M enumerated models).


% ═══════════════════════════════════════════════════════════════
\section{Discussion}
\label{sec:discussion}
% ═══════════════════════════════════════════════════════════════

\subsection{Relationship to the theoretical proposals}

The implementation captures the essential algorithmic content
of the split-forest approach~\cite{SiblingCompleted,CoverTree}
but omits several components of the full theory:

\begin{enumerate}[label=(\roman*)]
  \item \textbf{Rank-$d$ signatures and explicit blocking.}
    The type-set search implicitly achieves the effect of blocking
    (the search space is finite because $|\Tp(C_0)|$ is finite).
    The full rank-$d$ signature hierarchy, which includes parent
    summaries, child multiplicities, interface descriptors, and
    witness menus, is not constructed.

  \item \textbf{Descriptor language and local coherence axioms
    (C1--C5).}
    The singleton-composition check captures the most critical
    coherence constraint (when $\comp(\inv(R_1), R_2)$ is a
    singleton, the relation between witnesses is forced).
    The full C3--C4 downward-support axioms, which constrain how
    cross-edge domains evolve through ancestor-descendant chains,
    are not explicitly checked.  No test failures have been
    observed, but this is empirical evidence, not a proof of
    correctness.

  \item \textbf{Explicit model construction.}
    The implementation determines SAT/UNSAT but does not construct
    an explicit $\ALCIRCC{5}$ model.  The theoretical argument
    constructs a model via canonical refinements of finite-prefix
    disjunctive networks (K\"onig's lemma), then quotients by
    $\EQ$.

  \item \textbf{Completeness direction.}
    The theoretical proof that every satisfiable concept has a
    valid finite rank-$d$ quotient (the model-to-quotient
    extraction) is ``condensed'' in~\cite{SiblingCompleted}.
    Our implementation's correctness for the completeness direction
    is supported by the 902-test cross-validation but not by a
    formal proof.
\end{enumerate}

\subsection{Why the two reasoners agree}

The quasimodel reasoner and the cover-tree tableau use different
constraint-checking strategies but agree on all 911~tests.
This is because:
\begin{itemize}[nosep]
  \item Both start from the same closure, types, and $\Safe$
    relations.
  \item Both require demand closure (every $\exists R.D$ has a
    witness).
  \item The quasimodel reasoner's path-consistency check on the
    full 4-relation $\Safe$ network is \emph{strictly stronger}
    than the cover-tree tableau's $\{\DR,\PO\}$-restricted check.
    Thus any type set accepted by the quasimodel reasoner is
    also accepted by the cover-tree tableau (but not vice versa).
  \item The quasimodel reasoner's sibling check on all demand
    pairs is \emph{strictly stronger} than the cover-tree's
    cross-only check.
  \item The cover-tree's singleton-composition check catches
    exactly the cases where the quasimodel's stronger checks
    would also reject (all singleton-composition UNSAT patterns
    are caught by both).
\end{itemize}

In principle, the cover-tree tableau could accept a type set that
the quasimodel reasoner rejects (since its checks are weaker in
some dimensions).  The fact that no such case arose in 911~tests
suggests that the weaker checks are nonetheless sufficient, which
is consistent with the theoretical claim of the split-forest
approach.

\subsection{Open questions}

\begin{enumerate}[label=(\arabic*)]
  \item \textbf{Formal correctness.}
    Does the cover-tree tableau as implemented constitute a sound
    and complete decision procedure?  Soundness requires showing
    that every accepted type set can be unfolded into a genuine
    $\ALCIRCC{5}$ model.  Completeness requires showing that every
    satisfiable concept has an accepted type set.  Both directions
    are supported by the 902-test cross-validation and by the
    theoretical arguments of~\cite{SiblingCompleted}, but a
    complete formal proof is outstanding.

  \item \textbf{Complexity.}
    The implementation's worst-case complexity is not analyzed.
    The type-set search is bounded by $2^{|\Tp(C_0)|}$, which
    is doubly exponential in $|\cl(C_0)|$.  Whether the
    cover-tree approach yields an $\EXPTIME$ decision procedure
    (matching the known lower bound for $\ALCI$) is open.

  \item \textbf{Extension to RCC8.}
    RCC5 has the patchwork property; RCC8 does not.  The
    split-forest approach relies on full tractability of
    disjunctive RCC5 networks.  Extending the approach to
    $\ALCIRCC{8}$ would require a different strategy for the
    cross-edge network.
\end{enumerate}


% ═══════════════════════════════════════════════════════════════
\section{Conclusion}
\label{sec:conclusion}
% ═══════════════════════════════════════════════════════════════

We have implemented a cover-tree tableau decision procedure for
$\ALCIRCC{5}$ concept satisfiability, based on the split-forest /
sibling-interface approach.  The key algorithmic contributions are:
\begin{enumerate}[nosep]
  \item Separation of demands into tree ($\PP/\PPI$) and cross
    ($\DR/\PO$) categories, with distinct constraint-checking
    strategies for each.
  \item A restricted sibling compatibility check that only
    considers $\DR/\PO$ witness pairs, avoiding problematic
    mixed $\PP/\DR$ compositions.
  \item A singleton-composition propagation check that catches
    cross-role UNSAT patterns where the forced relation between
    a tree witness and a cross witness violates type-safety.
\end{enumerate}

The implementation passes 911~test concepts with zero mismatches
against an independent quasimodel-based reasoner.  Independent model
verification confirms 678/775~SAT concepts with concrete models
(zero verification failures), and a cover-tree decomposition test
shows that 89.0\% of built models naturally exhibit cover-tree
structure, with \textbf{100\%} of SAT concepts (775/775) having finite
cover-tree models and zero genuine counterexamples across
12.2\,M enumerated models (Section~\ref{sec:decomposition}).
Together, these three
layers provide strong computational evidence for the correctness of
the split-forest approach.  A formal correctness proof, complexity
analysis, and extension to $\ALCIRCC{8}$ remain open.


% ═══════════════════════════════════════════════════════════════
% Bibliography
% ═══════════════════════════════════════════════════════════════

\begin{thebibliography}{10}

\bibitem{Wessel2003}
Michael Wessel.
\newblock Qualitative spatial reasoning with the $\mathcal{ALCI}_{RCC}$
  family---first results and unanswered questions.
\newblock Technical Memo No.~324/03, University of Hamburg, 2003.

\bibitem{Wessel2002}
Michael Wessel.
\newblock On spatial reasoning with description logics---position paper.
\newblock In \emph{Proc.\ International Workshop on Description Logics
  (DL~2002)}, CEUR-WS, 2002.

\bibitem{ClaudeDecidability}
Claude (Anthropic), prompted by Michael Wessel.
\newblock On the decidability of $\mathcal{ALCI}_{RCC5}$ and
  $\mathcal{ALCI}_{RCC8}$: Quasimodels meet the patchwork property.
\newblock Manuscript, March 2026 (revised).

\bibitem{GPTContextual}
GPT-5.4 Pro (OpenAI), prompted by Michael Wessel.
\newblock $\mathcal{ALCI}_{RCC5}$ contextual tableau draft.
\newblock Manuscript, April 2026.

\bibitem{SiblingCompleted}
GPT-5.4 Pro (OpenAI), prompted by Michael Wessel.
\newblock Finite sibling-interface descriptors and a proposed completion
  for $\mathcal{ALCI}_{RCC5}$: completed split-forest version with
  canonical refinement and typed-$\EQ$ quotienting.
\newblock Manuscript, April 2026.

\bibitem{CoverTree}
GPT-5.4 Pro (OpenAI), prompted by Michael Wessel.
\newblock A proposed cover-tree tableau calculus for
  $\mathcal{ALCI}_{RCC5}$ with eventuality-aware signature blocking.
\newblock Manuscript, April 2026.

\bibitem{RenzNebel1999}
Jochen Renz and Bernhard Nebel.
\newblock On the complexity of qualitative spatial reasoning:
  A maximal tractable fragment of the Region Connection Calculus.
\newblock \emph{Artificial Intelligence}, 108(1--2):69--123, 1999.

\end{thebibliography}

\end{document}
